% =============================================================================
\section{Simulation Methodology}
\label{app:simulation}
% =============================================================================

This appendix provides the complete simulation methodology, parameter tables,
and random seed generation procedure for the Monte Carlo simulations presented
in Section~\ref{sec:simulations}.

\subsection{Simulation Engine}
\label{subsec:sim-engine}

The simulation engine is implemented in Python 3.13 using NumPy 1.26 and
SciPy 1.12. Each simulation run models a 48-month (1{,}461-day) horizon with
daily time steps. The state vector at each step includes:

\begin{itemize}
  \item Circulating supply $S(t)$
  \item Total \veCORTEX{} supply $V_{\text{total}}(t)$
  \item Number of active agents $N(t)$
  \item Treasury balance $R_{\text{treasury}}(t)$
  \item Cumulative tokens burned $B(t)$
  \item Market regime $\mu(t) \in \{\text{Bull, Neutral, Bear}\}$
  \item Per-agent state: bonding curve position, graduation status, trading
    volume, composability connections
\end{itemize}

\subsection{Parameter Tables}
\label{subsec:sim-params}

Table~\ref{tab:sim-core} presents the core simulation parameters.

\begin{table}[H]
\centering
\caption{Core simulation parameters.}
\label{tab:sim-core}
\begin{tabular}{@{}lll@{}}
\toprule
\textbf{Parameter} & \textbf{Value} & \textbf{Description} \\
\midrule
$T_{\text{sim}}$ & 1{,}461 days & Simulation horizon (4 years) \\
$N_{\text{runs}}$ & 1{,}000 & Monte Carlo runs per scenario \\
$\Delta t$ & 1 day & Time step \\
$S_0$ & $5.75 \times 10^{14}$ & Initial circulating supply \\
$R_{\text{total}}$ & $2.5 \times 10^{15}$ & Total staking emissions \\
$\mathbf{w}$ & $(0.4, 0.3, 0.2, 0.1)$ & Annual emission weights \\
$f_{\text{base}}$ & 0.005 & Base fee rate \\
$f_{\text{dyn}}$ & 0.015 & Dynamic fee coefficient \\
$\sigma_{\max}$ & 0.5 & Maximum volatility for fee scaling \\
$G$ & 50{,}000 CORTEX & Graduation threshold \\
$T_{\max}$ & 4 years & Maximum lock duration \\
\bottomrule
\end{tabular}
\end{table}

Table~\ref{tab:sim-growth} presents the growth scenario parameters.

\begin{table}[H]
\centering
\caption{Growth scenario parameters.}
\label{tab:sim-growth}
\begin{tabular}{@{}lccc@{}}
\toprule
\textbf{Parameter} & \textbf{Conservative} & \textbf{Moderate} & \textbf{Aggressive} \\
\midrule
Agent arrival rate $\lambda_0$ & 1.5/week & 3/week & 6/week \\
Growth rate $\alpha$ & 0.01 & 0.02 & 0.04 \\
Mean daily volume (log) $\mu_V$ & 7.0 & 8.0 & 9.0 \\
Volume std dev $\sigma_V$ & 1.5 & 1.5 & 1.5 \\
Graduation rate & 40\% & 60\% & 80\% \\
Staking participation & 25\% & 40\% & 55\% \\
\bottomrule
\end{tabular}
\end{table}

Table~\ref{tab:sim-regime} presents the market regime transition matrix.

\begin{table}[H]
\centering
\caption{Daily market regime transition probabilities.}
\label{tab:sim-regime}
\begin{tabular}{@{}lccc@{}}
\toprule
\textbf{From $\backslash$ To} & \textbf{Bull} & \textbf{Neutral} & \textbf{Bear} \\
\midrule
Bull & 0.97 & 0.02 & 0.01 \\
Neutral & 0.03 & 0.94 & 0.03 \\
Bear & 0.01 & 0.04 & 0.95 \\
\bottomrule
\end{tabular}
\end{table}

The transition probabilities are calibrated to produce regime durations
consistent with historical crypto market cycles: bull regimes last
approximately 33 days on average ($1/0.03$), neutral regimes approximately
17 days ($1/0.06$), and bear regimes approximately 20 days ($1/0.05$).

Table~\ref{tab:sim-staker} presents the staker behavior parameters.

\begin{table}[H]
\centering
\caption{Staker behavior distribution.}
\label{tab:sim-staker}
\begin{tabular}{@{}lcc@{}}
\toprule
\textbf{Lock Duration} & \textbf{Probability} & \textbf{veCORTEX Multiplier} \\
\midrule
4 years ($T_{\max}$) & 30\% & 1.0 \\
2 years ($T_{\max}/2$) & 40\% & 0.5 \\
1 year ($T_{\max}/4$) & 30\% & 0.25 \\
\bottomrule
\end{tabular}
\end{table}

\subsection{Volatility Model}
\label{subsec:sim-vol}

Realized volatility is modeled as a regime-dependent log-normal process:

\begin{equation}
  \sigma(t) \sim \text{LogNormal}(\mu_\sigma(\mu(t)),\; \xi^2)
\end{equation}

with regime-dependent mean parameters:

\begin{table}[H]
\centering
\caption{Volatility model parameters by regime.}
\label{tab:sim-vol}
\begin{tabular}{@{}lcc@{}}
\toprule
\textbf{Regime} & $\mu_\sigma$ (log-scale) & \textbf{Implied Mean Vol} \\
\midrule
Bull & $-1.2$ & 33\% (annualized) \\
Neutral & $-1.8$ & 18\% \\
Bear & $-0.8$ & 49\% \\
\bottomrule
\end{tabular}
\end{table}

The shared variance parameter is $\xi^2 = 0.3$, producing fat-tailed
volatility distributions consistent with empirical crypto market data.

\subsection{Trading Volume Model}
\label{subsec:sim-volume}

Daily trading volume per agent follows:

\begin{equation}
  V_{\text{agent}}(t) = \exp(\mu_V + \sigma_V Z) \cdot \kappa(\mu(t)) \cdot \text{age\_factor}(t - t_{\text{create}})
\end{equation}

where $Z \sim \mathcal{N}(0,1)$, $\kappa(\text{Bull}) = 1.5$,
$\kappa(\text{Neutral}) = 1.0$, $\kappa(\text{Bear}) = 0.5$, and:

\begin{equation}
  \text{age\_factor}(\Delta) = \begin{cases}
    0.3 + 0.7 \cdot \Delta / 30 & \text{if } \Delta < 30 \\
    1.0 & \text{if } 30 \leq \Delta < 180 \\
    1.0 \cdot e^{-0.002(\Delta - 180)} & \text{if } \Delta \geq 180
  \end{cases}
\end{equation}

This models initial ramp-up, peak activity, and gradual decay.

\subsection{Composability Model}
\label{subsec:sim-comp}

Agent tool registration follows a Bernoulli process with probability
$p_{\text{tool}} = 0.7$ (70\% of agents expose at least one tool). Cross-agent
invocations follow a preferential attachment model:

\begin{equation}
  P(\text{agent } j \text{ calls agent } i) \propto CS_i^{0.8} + \epsilon
\end{equation}

where $\epsilon = 0.01$ ensures nonzero probability for agents with low
Composability Scores, preventing rich-get-richer lock-in.

\subsection{Governance Model}
\label{subsec:sim-gov}

The governance simulation generates Senator proposals at a rate of 2 per month,
with:
\begin{itemize}
  \item 60\% Parameter proposals, 25\% Treasury proposals, 10\% Constitution
    proposals, 5\% Emergency proposals.
  \item Voter turnout drawn from a beta distribution
    $\text{Beta}(2, 3)$ (mean 40\%, right-skewed).
  \item Delegation probability: 35\% of stakers delegate, with delegates
    chosen proportional to \veCORTEX{} weight.
  \item Vote correlation with Senator recommendation: 0.7 (voters are
    partially influenced by the Senator's proposal quality assessment).
\end{itemize}

\subsection{Random Seed Generation}
\label{subsec:seeds}

Each of the 1{,}000 Monte Carlo runs uses a distinct random seed:

\begin{equation}
  \text{seed}_k = \text{SHA-256}(\text{master\_seed} \,\|\, k) \mod 2^{32}
\end{equation}

where $\text{master\_seed} = \texttt{0xDEADBEEF42}$ and $k \in \{0, 1,
\ldots, 999\}$. The SHA-256 hash ensures statistical independence between
runs. The master seed is fixed for reproducibility; all results can be exactly
replicated by using the same master seed and simulation code.

\subsection{Stress Test Methodology}
\label{subsec:sim-stress-method}

The stress test (Figure~\ref{fig:stress-test}) is implemented as a
deterministic shock at day 365 (month 12):

\begin{itemize}
  \item Token price: instantaneous 60\% decline.
  \item Trading volume: 40\% decline (partially offset by panic selling).
  \item New agent arrivals: 80\% decline for 30 days, linear recovery over 90
    days.
  \item Staker behavior: 20\% of non-locked stakers exit (locked stakers
    cannot exit by construction).
\end{itemize}

Recovery is measured as the number of days until TVL returns to the pre-crash
level.

\subsection{Confidence Interval Construction}
\label{subsec:sim-ci}

Confidence intervals in Figures~\ref{fig:supply-trajectory}--\ref{fig:sensitivity}
are constructed as pointwise 95\% intervals:

\begin{equation}
  \text{CI}_{95}(t) = \left[\hat{Q}_{0.025}(t),\; \hat{Q}_{0.975}(t)\right]
\end{equation}

where $\hat{Q}_p(t)$ is the empirical $p$-quantile across the 1{,}000 runs at
time $t$. We use quantile-based intervals rather than mean $\pm$ 1.96 standard
deviations because several output distributions (particularly trading volume
and burn amounts) are right-skewed.
